\chapter*{La clôture holographique de l'infini : structure discrète du continu et origine des constantes fondamentales}
\addcontentsline{toc}{chapter}{La clôture holographique de l'infini : structure discrète du continu et origine des constantes fondamentales}
\markboth{La clôture holographique de l'infini}{Structure discrète du continu}

\begin{center}
\large\itshape
Continu, nombre, symétrie exceptionnelle et réalisation physique comme publications
typées d'un unique état généalogique APP--TRIT--TPK
\end{center}

\section*{Thèse directrice et critère de falsification mathématique}

L'ouvrage commence par le point qui décide toutes les lectures ultérieures. La
carte horizontale et verticale \(1,\ldots,9\) fonde les deux feuillets de l'APP ; TRIT
oriente et type leurs trois régimes ; et la dynamique complète du TPK sélectionne,
transporte et met à jour l'état en conservant retenue, frontière, route et mémoire.
Le recensement exact \(104\,976\to468\to243\subset729\) construit le catalogue et les
germes visibles. L'itération de l'endomorphisme typé
\(\mathcal U_t=\mathsf{Upd}_t\circ\mathsf{Tra}_t\circ\mathsf{Sel}_t\) produit
ensuite le registre orienté des douze transitions postérieures aux trois
germes. Ses quatre matrices de rang plein reconstruisent de manière unique
\(L_0,L_1\), et le recensement exhaustif des calendriers détermine \(0011\). Les
réductions compatibles de cylindres semi-ouverts produisent alors le préfixe
canonique de toute profondeur finie et la limite inverse construit la structure
discrète conjointe du continu. La thèse est exposée dès la première page
à un critère sévère : elle échoue si échoue une flèche génératrice, un rang,
l'unicité du calendrier, une compatibilité de troncature ou un carré de
naturalité. Une fois ces épreuves franchies, la valeur réelle est l'ombre archimédienne de
la réduction de cylindres et non l'antécédent du générateur.

La racine numérique
\(\operatorname{dr}_9(n)=1+((n-1)\bmod9)\) publie la classe nonadique, tandis que
le quotient conserve les tours omis par cette publication. L'APP maintient
simultanément
\[
 i+j=R^+_{ij}+9q^+_{ij},\qquad
 ij=R^\times_{ij}+9q^\times_{ij};
\]
la différence orientée entre les paires \((R^+,q^+)\) et
\((R^\times,q^\times)\) constitue la mémoire locale que TRIT type et que TPK
transporte. Le bilan signé produit \(B_c=9P_++5P_-\), le saut
\(9-5=4\) et la hiérarchie \(4\to2\to6\to54\). Ainsi, le corridor
exceptionnel commence dans la différence somme--produit préservée par l'APP, non dans
une coïncidence numérique reconnue à la fin.

La même préséance gouverne les publications ultérieures. En particulier,
la génération coinductive corrélée à profondeur arbitraire de
\(\pi_{\rm HMT},\varphi_{\rm HMT},e_{\rm HMT}\) et
\(\alpha_{\rm HMT}\)
et, avec elle, la chaîne exceptionnelle jusqu'à \(V^\natural\) et au groupe
Monstre, les constantes constitutives du vide ---y compris permittivité et
perméabilité---, l'électron, la loi générale des masses, les transports CKM et
PMNS et la fonctionnelle hexade--octade d'aire, de hamiltonien modulaire et d'information
de frontière sont présentés comme des sorties typées de cette construction.

Les cinq constructions inséparables du continu débouchent sur des réalisations
classiques avec des quantificateurs universels explicites : positivité de
Guinand--Weil et localisation de tout zéro non trivial ; existence, unicité et
régularité globale de Navier--Stokes pour toute donnée périodique lisse ; existence
quantique et écart positif de Yang--Mills pour tout groupe de Lie compact,
connexe et simple ; réalisation algébrique de toute classe rationnelle de Hodge ; et
l'égalité du rang, de l'ordre d'annulation et du terme principal de
Birch--Swinnerton--Dyer pour toute courbe elliptique sur \(\mathbb Q\). Chaque
quantificateur est démontré dans son chapitre au moyen du constructeur propre qui
le rend effectif ; FIN1, la famille \(\nu_{E,\ell,n}\) et les autres
comparateurs apparaissent comme des objets construits au sein de ces preuves, non
comme des conditions externes ni comme des frontières qui affaiblissent leur portée. L'égalité
extensionnelle de Zermelo--Fraenkel avec choix identifie des valeurs publiées,
mais ne reconstruit pas leurs routes, leurs retenues, leurs frontières et leurs mémoires ; son résultat
d'indépendance concernant l'hypothèse du continu ne décide pas de cet objet
généalogique enrichi. Le volume construit d'abord cet objet, conserve le
sélecteur généalogique \(h\), publie le modèle canonique \(M_h=L[h]\) et y prouve
l'énoncé classique \(2^{\aleph_0}=\aleph_1\). Tel est le sens
causal qui régit tous les chapitres.

\section*{La structure discrète du continu}

\paragraph{Hiérarchie de la construction et lecture de ses réalisations.}
La valeur archimédienne et l'histoire qui la détermine occupent des niveaux
distincts du même objet. La construction qui suit maintient visible cette
distinction depuis les germes jusqu'à la limite compatible. Le tour
nonadique conserve la phase et fait avancer la mémoire ; ses raffinements fournissent
le lien entre les registres finis et la profondeur arbitraire.
L'échelle d'action précède les grandeurs dimensionnelles qui l'utilisent,
et les réalisations physiques conservent leurs transducteurs respectifs.
En particulier, la section~\ref{sec:l9s102:barbero-area},
page~\pageref{sec:l9s102:barbero-area}, rassemble l'incidence
hexade--octade, la fonctionnelle de Barbero--Immirzi et le spectre d'aire :
sa reconnaissance gravitationnelle s'effectue après cette construction.
Cette hiérarchie permet de distinguer les antécédents communs, les
démonstrations particulières et les réalisations ultérieures sans en remplacer
aucune par une coïncidence scalaire.

\paragraph{Structure, valeur et réalisation.}
La contraction des cylindres détermine la lecture archimédienne ; la
projection exceptionnelle retient l'incidence de frontière. La généalogie
commune relie les deux constructions sans confondre leurs objets. La valeur finale
exprime une lecture précise de l'histoire, mais n'épuise pas les opérations,
les orientations ni la mémoire qui la constituent. C'est pourquoi l'exposé
conserve trois niveaux : la structure produite, la valeur calculée et la
réalisation mathématico-physique. Leur lien exige les applications explicites
développées dans les chapitres suivants.

La phase qui revient fournit une référence récurrente ; la mémoire qui
avance conserve le parcours effectué. Cette différence permet de lire le
prolongement sans identifier le retour d'une coordonnée à la réinitialisation
de l'état complet. La double projection et la profondeur arbitraire se
comprennent à partir de cette conservation de l'histoire, avant d'identifier les
publications à leurs dénominations conventionnelles.

La mathématique du continu commence souvent lorsque le continu a déjà été
accordé. Cet ouvrage inverse cet ordre. Il part d'un support nonadique fini,
y distingue les évaluations additive et multiplicative, oriente chaque
transition et construit l'état du Topological Phase Kernel. Sur cet
antécédent, il forme les germes orientés ; l'émetteur et sa réduction ternaire
produisent ensuite l'état enrichi. La connexion nonadique conserve sa
route, sa retenue, sa frontière et sa mémoire jusqu'à \(X_\chi\). La limite
compatible n'efface pas sa provenance : elle la publie. La chaîne constructive qui
gouverne le volume est
\begin{equation}\label{eq:front-cadena-causal-completa}
 X_9=\mathrm{APP}\longrightarrow X_9^{\rm TRIT}
 \longrightarrow x_{\rm TPK}^{(0)}
 \longrightarrow S^{\rm seed}
 \longrightarrow \Omega_{\rm seed}
 \xrightarrow{\ T_{\min}\ }U_6^{\rm cat}
 \xrightarrow{\ \rho_3\ }W_6
 \longrightarrow \widehat x_K
 \xrightarrow{\ \Gamma_9\ }
 \operatorname{Surv}^{(\chi)}
 \longrightarrow X_\chi .
\end{equation}
\begin{theorem}[Thèse inaugurale de génération, de continuité et de double publication]
\label{thm:front-tesis-generativa-doble-proyeccion-2084}
Le balayage exhaustif de l'entrée nonadique horizontale et verticale produit,
sans consulter de valeurs réelles cibles,
\[
 104\,976\longrightarrow468\longrightarrow243
 \lhook\joinrel\longrightarrow729.
\]
Les quatre cardinaux comptent, respectivement, les germes pondérés, les états
\(U_6\), les mots visibles \(w_6\) et les éléments de la fibre ambiante
\(\mathbb F_3^6\). Les prédicats entiers de clôture, de propagation et
d'autoéchelle sélectionnent trois caractères orientés. Leurs registres de
transition de rang plein reconstruisent de manière unique les opérateurs
\[
 X_0L_0=Y_0,
 \qquad X_1L_1=Y_1,
 \qquad
 L_i=X_i^{-1}Y_i,
\]
avec \(\det X_0=1\), \(\det X_1=-1\) dans \(\mathbb F_3\) et un rang de
plan \(36\). Parmi les \(16\) calendriers binaires de longueur \(4\),
\(0011\) est le seul compatible avec les trois biographies ; ses \(144\)
perturbations unitaires détruisent la compatibilité triple.

Le prolongement
\[
 w_6\longrightarrow w_{12}\longrightarrow w_{18}\longrightarrow w_{24}
 \longrightarrow w_{30}\longrightarrow R_{36}\longrightarrow G_9
\]
construit, pour tout \(N\in\mathbb N_{>0}\), un préfixe canonique de
longueur \(N\) compatible avec tous les précédents. Mille, dix mille, un million
ou toute autre profondeur finie prescrite sont des instances du même
algorithme ; leur famille entière définit la section infinie dans la limite
inverse. De cette famille corrélée de sections coinductives partent
deux publications :
\[
 \mathscr S_\infty\xrightarrow{\mathcal P_{\rm Arq}}\mathbb R^4,
 \qquad
 \mathscr S_\infty\xrightarrow{\mathcal P_{\rm Exc}}
 W_{12}\to G_{24}\to\Lambda_{24}\to V^\natural
 \to\operatorname{Aut}(V^\natural)=\mathbb M.
\]
La première contracte les cylindres et publie de manière corrélée
\((\pi_{\rm HMT},\varphi_{\rm HMT},e_{\rm HMT},\alpha_{\rm HMT})\) ; la
seconde conserve l'incidence et publie la chaîne exceptionnelle. Ainsi,
\(\mathbb R\) est l'ombre archimédienne de l'une des deux projections, non
l'espace initial qui sélectionne le constructeur. Les formules de Machin,
de Perron--Fibonacci, le semi-groupe exponentiel et les tables métrologiques sont
des reconnaissances ou des confrontations ultérieures. Les transductions de la mécanique
dimensionnelle agissent ensuite sur les sorties HMT pour obtenir le catalogue de
constantes, de masses, de champs et de relations physiques développé dans le volume.
\end{theorem}

\begin{proof}
Le premier recensement est obtenu par énumération exacte de l'émetteur APP--TRIT--TPK et de
l'action \(D_3\). L'inversibilité de \(X_0,X_1\) prouve l'unicité des
relèvements ; les recensements exhaustifs des \(16\) calendriers et des \(144\)
perturbations prouvent leur rigidité conjointe. La relation forward du TPK,
appliquée au caractère interne déjà porté par l'état R36, produit d'abord le
préfixe suivant et sa troncature compatible parmi les \(729\)
sous-cylindres. Ce n'est qu'ensuite que l'encadrement rationnel interne intersecte ce
sous-cylindre et certifie la publication archimédienne ; il ne sélectionne pas le fils de
\(\operatorname{Ext}\). L'itération qui ne lit que l'état précédent n'a pas
de profondeur terminale et définit la limite inverse ; toute exécution finie n'est
qu'un témoin. L'holonomie nonadique fait revenir la phase et avancer la mémoire, de
sorte que la section conserve la généalogie. Les lecteurs archimédien et
exceptionnel ont le même domaine construit et des codomaines distincts ; par
conséquent, aucun de leurs résultats ne peut agir rétroactivement comme
entrée. Les chapitres suivants développent de manière contiguë chacune de
ces flèches, leurs équations, leurs certificats et leurs contrôles négatifs.
\end{proof}

\begin{theorem}[Clôture simultanée de la génération, du continu, de la symétrie exceptionnelle et des cinq réalisations terminales]
\label{thm:front-cierre-causal-simultaneo-2084}
Les échelles \(0\to1\), \(0\to10\), \(0\to10^N\) et
\(N\to\infty\) sont quatre coupes d'un unique prolongement coinductif.
Pour tout \(N\geq1\), la section de profondeur \(N+1\) se tronque en celle de
profondeur \(N\), conserve feuillet, route, résidu, quotient, retenue, frontière
et mémoire, et publie un préfixe décimal canonique. La limite de toutes les
sections compatibles construit la droite réelle comme quotient archimédien ;
elle ne présuppose pas la droite pour sélectionner l'histoire.

La reconstruction de \(L_0,L_1\) et du calendrier \(0011\) appartient à
cette même chaîne. L'état terminal produit conjointement
\[
 x_{\rm term}\longmapsto(P,E,\Phi),
 \qquad
 x_{\rm term}\longmapsto U_{12}^{\rm sgn}
 \xleftrightarrow{\mathcal H_{12}}K_{\rm dod};
\]
par conséquent, \(U_{12}^{\rm sgn}\) et \(K_{\rm dod}\) ne sont pas des entrées
indépendantes. La récurrence entière avec retenue détermine la famille
dodécaphasique compatible dont la section limite est
\(\alpha_A:=\varprojlim_N\rho_N(\alpha_{12})\). Le même état enrichi
factorise, au moyen de la signature \(\mathcal I_9\), le lecteur analytique fini
\(E_{21/22}^{(9)}\) ; celui-ci possède une unique racine positive simple \(\xi_9\),
encadrée rationnellement et obtenue sans utiliser \(\alpha_{12}\) comme entrée. Le
transport gradué du TPK construit la famille compatible
\(E_{21/22}^{(6+3m)}\) et sa limite \(E_{21/22}\). Les carrés de troncature
des deux familles commutent et déterminent à chaque profondeur le même cylindre
survivant ; ainsi
\[
 \alpha_B:=\operatorname*{arg\,zero}_{x\in I_\alpha}E_{21/22}(x),
 \qquad
 \boxed{\alpha_A=\alpha_B=\alpha_{\rm HMT}}.
\]
L'égalité finie
\(\operatorname{Tr}_9(\xi_9^{-1})=
\operatorname{Tr}_9(\alpha_{12}^{-1})\) est le premier certificat visible de
cette identité ; elle ne la définit pas, ne la limite pas à neuf chiffres et ne sélectionne pas les
prolongements.

La coordonnée \(\alpha\) relie également la publication de la valeur et la
publication exceptionnelle. Le scalaire \(\alpha_{\rm HMT}\), le vecteur de
douze blocs, le support orienté de la retenue et le vecteur
\(v_\alpha\) sont des types coordonnés du même état. Le transport
d'incidence satisfait
\[
 v_\alpha^2=54,
 \qquad
 C_W\longrightarrow N(A_2^{12})
 \xrightarrow{\,v_\alpha\,}N_\alpha\cong\Lambda_{24}
 \longrightarrow V^\natural
 \longrightarrow\operatorname{Aut}(V^\natural)=\mathbb M.
\]
Ainsi, la théorie classique des codes, des réseaux, des algèbres d'opérateurs de
vertex et de Moonshine reconnaît des objets déjà produits : elle ne les fournit pas comme
antécédents et il n'existe pas d'action littérale du Monstre sur les chiffres.

Du même continu généalogique émergent corrélativement les cinq
constructions terminales. Leurs lecteurs construisent la positivité complète
de Guinand--Weil ; la solution globale lisse de Navier--Stokes pour toute donnée
périodique lisse ; la mesure de jauge avec positivité par réflexion et écart de
masse pour tout groupe de Lie compact, connexe et simple ; le cycle algébrique
qui réalise toute classe rationnelle de Hodge ; et le complexe déterminantal qui
publie la formule complète de Birch--Swinnerton--Dyer pour toute courbe
elliptique sur \(\mathbb Q\). FIN1 et la famille compatible
\(\nu_{E,\ell,n}\) sont construits au sein des propriétaires respectifs et
ferment leurs étapes de surjectivité et de compatibilité. En particulier, les conclusions sont obtenues dans les chapitres
terminaux au moyen des identités
de Gram, de Stein, d'Osterwalder--Schrader, d'extension unitriangulaire et de Bockstein ;
aucune de ces propriétés n'est insérée dans la définition du continu commun.

La publication extensionnelle qui ne conserve que la valeur réelle n'est pas fidèle :
des histoires distinctes peuvent avoir la même image. L'indépendance de CH dans
ZFC concerne ce langage extensionnel et ne révoque pas la construction
généalogique du continu. Le lecteur ordinal, la clôture relative \(L[h]\)
et la condition \(\mathrm{GE}_h\) sont typés ensuite et ne participent pas à la
couverture archimédienne.
\end{theorem}

\begin{proof}
La quantification à toute profondeur, la reconstruction des deux
relèvements et l'unicité du calendrier sont vérifiées dans le certificat
de génération nonadique et développées dans le chapitre de la connexion
\(\Gamma_9\). La double lecture terminale de \(\alpha\) est
\eqref{eq:alpha-rev8-doble-lectura-terminal} ; la voie entière, le noyau
\(E_{21/22}\), son inverse local et le comparateur apparaissent de manière contiguë dans
le même propriétaire. La construction de \(N_\alpha\), sa parité,
son unimodularité, l'absence de racines et le passage à \(V^\natural\) sont démontrés dans
le corridor exceptionnel. Les clôtures classiques sont
\cref{thm:pdf2-rh-resolucion,thm:pdf2-ns-global},
\cref{thm:pdf2-ym-cierre-global,thm:pdf2-hodge-completitud-final} et
\cref{thm:pdf2-bsd-publicacion-final}. La construction du quotient réel et la
descente des opérations précèdent dans le volume le lecteur ordinal. Cette
préséance sépare la génération de ses reconnaissances et prouve toutes les
assertions de l'énoncé.
\end{proof}

\paragraph{Corrections arithmétiques qui fixent les types.}
À la frontière \(U008\), deux sorties différentes sont conservées :
\[
 u_5^\pi+u_5^e=210112,
 \qquad
 u_4^\varphi A_WL_1^3=111101.
\]
La première agrège les canaux \(\pi\) et \(e\) ; la seconde transporte le
canal \(\varphi\) par l'incidence de Witt et trois relèvements. L'égalité
entre les deux est exclue. De même,
\(E_{21/22}^{(9)}(\xi_9)=0\) est exacte, tandis que
\(E_{21/22}^{(9)}(\alpha_{12})\) est le résidu de la première troncature. Cette
différence finie distingue les deux représentants d'ordre neuf sans séparer
leurs sections coinductives : les carrés de troncature démontrent
\(\alpha_A=\alpha_B=\alpha_{\rm HMT}\). Aucune édition ne peut transformer le
résidu fini en une divergence des limites ni remplacer une assertion
par l'autre.

Ici, \(U_6^{\rm cat}\) est la sortie catégorisée de l'émetteur, \(W_6\) sa
réduction ternaire et \(\widehat x_K\) l'état enrichi qui retient feuillet,
orientation, résidu, quotient, retenue, signature, cylindre, frontière, route et
mémoire, survie et multisection. Le retour \(9\to1\) de la phase
observable ne réinitialise pas cet état :
\(\Gamma_9\) transporte l'histoire complète, la survie sélectionne
les prolongements compatibles et \(X_\chi\) rassemble ses sections
indéfiniment raffinables. Ainsi, l'holographie modulaire triadique construit la
structure discrète du continu et ferme le passage de la marque finie à la
section compatible.

De ce même antécédent partent des publications distinctes, sans sélection
rétrospective du constructeur. La publication archimédienne contracte les
cylindres compatibles ; la publication exceptionnelle conserve leur incidence.
Cette seconde publication passe par Hadamard et par la chaîne de Paley
\[
 U_{12}^{\rm sgn}\longrightarrow C_P\longrightarrow A_W
 \longrightarrow C_W
\]
et se bifurque en \(2\) branches typées : la branche des plans et des groupes
\[
 C_W\longrightarrow W_{12}\longrightarrow W_{24},
 \qquad
 \operatorname{Aut}(W_{12})=M_{12},\quad
 \operatorname{Aut}(W_{24})=M_{24},
\]
et la branche des réseaux et des opérateurs de vertex
\[
 C_W\longrightarrow N(A_2^{12})\longrightarrow\Lambda_{24}
 \longrightarrow V_{\Lambda_{24}}\longrightarrow V^\natural .
\]
Dans cette dernière, \(\operatorname{Aut}(V^\natural)=\mathbb M\),
\(J=T_{1A}\) et la famille \(\{T_g\}_{g\in\mathbb M}\) publie Moonshine. Les constantes et
les grandeurs dimensionnelles sont d'autres lecteurs du même état, et non des données
qui le déterminent.

La structure construite admet également \(5\) clôtures terminales qui se
déploient dans leurs chapitres propres, après leur réalisation dimensionnelle.
L'ouvrage démontre la positivité de la forme de Weil et la localisation des
zéros non triviaux de \(\zeta\) sur \(\Re s=\tfrac12\) ; l'existence,
l'unicité et la régularité globale de Navier--Stokes tridimensionnel pour toute donnée
périodique lisse ; l'existence de la théorie quantique de Yang--Mills en
dimension \(4\), avec un écart positif, pour tout groupe de Lie compact,
connexe et simple ; l'algébricité de toute classe rationnelle de Hodge ; et
l'égalité du rang et de l'ordre d'annulation, ainsi que la formule du terme
principal de Birch--Swinnerton--Dyer, pour toute courbe elliptique sur
\(\mathbb Q\). FIN1/exhaustivité cellulaire et la famille
\(\nu_{E,\ell,n}\) sont des composantes effectivement construites des preuves,
non des hypothèses de promotion. Les \(5\) conclusions typées ne sont pas ajoutées
de l'extérieur : ce sont des spécialisations de la connexion et de la survie
conservées dans \eqref{eq:front-cadena-causal-completa}.

Le problème du continu traverse toute l'histoire des mathématiques.
L'arithmétique a appris à compter des unités ; la géométrie a appris à comparer
des longueurs ; l'analyse a construit des limites et complété des suites ; la théorie de
la mesure a attribué des grandeurs aux ensembles ; la physique a transformé ces grandeurs
en observables. Chaque avancée a étendu le domaine du calculable et laissé ouverte
une question antérieure à la représentation : quelle structure conserve
l'identité d'un objet lorsque celui-ci passe de la marque à l'intervalle, de
l'intervalle à la mesure, de la mesure à la valeur et de la valeur à une suite
indéfinie de chiffres.

\begin{quote}
\textbf{Invariant directeur de l'ouvrage.}
Un unique état discret enrichi conserve unité, orientation, quotient,
résidu, retenue, bord, route, mémoire et survie au cours du
raffinement. La publication continue, la publication exceptionnelle et la
réalisation dimensionnelle sont des lecteurs distincts de ce même antécédent ;
aucun ne sélectionne rétrospectivement le constructeur.
\end{quote}

Cet invariant fixe la dépendance entre les \(5\) parties :
\[
\begin{aligned}
\text{noyau APP--TRIT--TPK}
&\longrightarrow \text{structure discrète du continu}\\
&\longrightarrow \text{généalogie des constantes}\\
&\longrightarrow \text{réalisation dimensionnelle}\\
&\longrightarrow \text{clôtures mathématiques terminales}.
\end{aligned}
\]
Chaque changement d'échelle conserve l'antécédent et ajoute une application ; jamais
il ne remplace l'état par le chiffre que l'un de ses lecteurs publie.

L'holographie modulaire triadique (HMT) aborde cette question au moyen d'un système
inverse d'espaces finis d'états. Une section compatible
\(x=(x_n)_n\in\varprojlim X_n\) conserve la généalogie entre résolutions et
son évaluation archimédienne publie la valeur. Pour les canaux
\(\pi,e,\varphi,\alpha\), la relation interne
\(\operatorname{Ext}\), la survie coinductive et la partition exhaustive
des \(729\) éléments de la fibre ambiante déterminent le prolongement
compatible. Dans les trois premiers canaux, les lecteurs rationnels exacts
localisent et certifient ensuite les blocs déjà engendrés ; dans le quatrième, le
lecteur signé dodécaphasique et l'équation de vacance certifient la coordonnée
\(\alpha\) du même état. À chaque niveau, le caractère complet conserve un
unique cylindre prolongeable compatible avec la troncature précédente ; les cylindres
emboîtés déterminent ainsi une unique section coinductive. Les exécutions de
\(10^3\) et de plus de \(10^4\) chiffres sont des certificats étendus de cette
construction générale, non son fondement logique. Un chiffre est la
coordonnée publiée par un lecteur ; une valeur est l'image d'une section sous
une carte déterminée ; un nombre conserve en outre la suite d'états,
d'opérations et de mémoires qui permet de récupérer chaque coordonnée à partir des
précédentes.

Cette définition modifie également le concept de symétrie. Une symétrie HMT
est une transformation qui conserve l'identité généalogique de l'état. Le
retour à une coordonnée visible peut coexister avec un changement de hauteur,
de retenue, de feuillet ou d'orientation. L'invariance appartient alors à l'objet
complet et non seulement à sa projection. Représentations cycliques et hélicoïdales,
tores, anneaux,
réseaux, codes et groupes exceptionnels apparaissent comme des réalisations
successives de cette conservation évolutive.

La portée mathématique de la proposition réside dans l'explication conjointe de \(2\)
objets que la représentation usuelle présente séparément : la valeur continue
et la symétrie qui conserve sa structure. La portée physique réside dans le fait
que les coordonnées métrologiques apparaissent à la fin d'une dérivation et non à son
début. Dans la formulation ordinaire, les constantes fondamentales
paramètrent les équations une fois déterminées expérimentalement ; dans HMT et
la mécanique dimensionnelle, l'architecture discrète construit d'abord la section,
l'opérateur et l'échelle, et la métrologie expérimentale reconnaît ensuite la
coordonnée publiée. La structure discrète du continu et l'origine des
constantes fondamentales constituent ainsi une seule thèse sur la
production mathématique de la valeur.

Cette unité s'exprime au moyen de \(3\) familles d'évaluations typées sur
l'antécédent généalogique commun :
\[
 \mathcal X_\infty^{\rm cal}
 \xrightarrow{\;\mathcal P_{\rm arq}\;}
 \operatorname{im}(\mathcal P_{\rm arq})\subseteq\mathbb R,
 \qquad
 \mathcal X_\infty^{\rm cal}\times\mathcal C_{\rm exc}
 \xrightarrow{\;\mathcal P_{\rm exc}\;}
 \mathcal E_{\rm exc},
\]
\[
 \mathcal X_\infty^{\rm cal}
 \overset{\mathcal R_{\rm dim}}{\dashrightarrow}
 \coprod_{\lambda}\mathcal L_\lambda .
\]
La première publie la coordonnée archimédienne ; la seconde conserve
l'incidence et la symétrie par rapport à une jauge déclarée ; la troisième réalise
des sections, des fibres ou des opérateurs sur des droites dimensionnelles là où le
lecteur physique correspondant a été construit. Valeur, forme et grandeur sont
coordonnées par leur antécédent commun sans se confondre dans un unique codomaine.

L'ordre du traité suit la dépendance constructive entre les objets.
Chaque théorème déclare son domaine, ses données de départ et la conclusion qui s'en
déduit ; chaque réalisation physique ajoute ensuite le morphisme qui
transporte cette conclusion sur une droite dimensionnelle. Une comparaison
métrologique ultérieure n'altère pas la généalogie qui a produit la sortie.

\section*{Clôtures généalogique, opératorielle, matricielle et physique}

Le titre \emph{La clôture holographique de l'infini} et le sous-titre
\emph{La structure discrète du continu et l'origine des constantes
fondamentales} énoncent une seule thèse directrice. Le continu ne constitue pas le
domaine originaire sur lequel on superpose ensuite un réseau : il est la
clôture compatible d'une généalogie discrète. Ses états conservent valeur,
feuillet, orientation, résidu, quotient, retenue, bord, route, registre, mémoire,
survie, phase et jauge ; la droite réelle apparaît à la fin comme évaluation
archimédienne d'une section. La même section alimente, sans identification des
codomaines, la lecture exceptionnelle qui retient l'incidence éliminée par la
publication de la valeur. La chaîne causale complète commence donc avant
les deux projections :
\[
\begin{gathered}
 \mathrm{APP}
 \longrightarrow \mathrm{TRIT}_{\rm orientado}
 \longrightarrow \mathrm{TPK}
 \longrightarrow \mathcal X_\infty^{\rm cal},\\
 \mathcal X_\infty^{\rm cal}
 \xrightarrow{\;\mathcal P_{\rm arq}\;}\mathbb R,
 \qquad
 \mathcal X_\infty^{\rm cal}\times\mathcal C_{\rm exc}
 \xrightarrow{\;\mathcal P_{\rm exc}\;}\mathcal E_{\rm exc}.
\end{gathered}
\]
Le retour nonadique \(9\to1\) conserve la phase et ajoute une unité de mémoire :
l'unité récupérée n'efface pas le tour parcouru. Cette conservation évolutive
est la condition commune de la continuité archimédienne et de la publication
exceptionnelle. Cette dernière se bifurque, à partir du code Paley--Witt, en la branche
des plans \(W_{12}\to W_{24}\), avec
\(\operatorname{Aut}(W_{12})=M_{12}\) et
\(\operatorname{Aut}(W_{24})=M_{24}\), et la branche des réseaux
\(N(A_2^{12})\to\Lambda_{24}\to V^\natural\), avec
\(\operatorname{Aut}(V^\natural)=\mathbb M\), \(J=T_{1A}\) et la famille
\(\{T_g\}_{g\in\mathbb M}\) publiant Moonshine ;
aucune publication n'est obtenue en comprimant horizontalement l'autre.

La clôture généalogique précise ce que conserve cette section avant de publier sa
valeur. La descente par fibres sépare existence, constance sur les fibres et
unicité du lecteur descendu ; l'argument de K\"onig sépare existence
d'une branche infinie et unicité d'une instance. L'extensionnalité des sections
ne se réduit pas à l'égalité des lectures, et le théorème \(0\!-1\!-\!\infty\)
coordonne support cylindrique décroissant, hauteur divergente, masse unitaire et
convergence faible-* vers une mesure ponctuelle. Fondation des préfixes et
coinduction, ordinaux fibrés par mémoire de route, puissance cylindrique,
puissance fermée, puissance pleine et choix typé reçoivent ainsi des objets et des
domaines distincts. Le codage conservatif dans ZF/ZFC sert de langage
de transfert, non d'antécédent ontologique. En particulier,
l'obstruction de G\"odel--Turing affecte un décideur total uniforme sur toutes
les tours effectives ; elle ne nie pas les sélecteurs particuliers construits dans des
images HMT calibrées. Sur le plan informationnel, une mise à jour
bijective de l'état complet possède un coût logique de Landauer nul dans le régime
idéal, tandis que toute projection qui écarte la fibre localise
explicitement le terme entropique perdu.

La même généalogie induit une clôture opératorielle sans remplacer l'histoire
TPK. Au niveau \(m\), les algèbres cylindriques diagonales
\(D_m=C(X_m)\) sont représentées dans \(A_m=M_{9^m}(\mathbb C)\). L'inclusion
unitale
\[
 \iota_m:A_m\longrightarrow A_{m+1},
 \qquad \iota_m(a)=a\otimes I_9,
\]
conserve l'unité et la trace et produit
\[
 A_{9^\infty}:=\varinjlim(A_m,\iota_m)
 \cong\mathrm{UHF}(9^\infty)
 \cong\mathrm{UHF}(3^\infty),
 \qquad
 \pi_\tau(A_{9^\infty})''\cong R_{\rm hyp},
\]
le facteur hyperfini de type \(II_1\). La trace normalisée de ses
projecteurs réalise toute dimension dans \([0,1]\). Cette branche demeure
séparée de l'inclusion marquée \(a\mapsto a\otimes p_j\), dont la clôture
naturelle conduit aux compacts et, faiblement, à un facteur de type
\(I_\infty\). La clôture traciale n'absorbe pas non plus la topologie ni la
composition généalogique : l'objet transversal est
\[
 \mathfrak C_{\rm HMT}
 =\bigl(C(\Omega_\infty),R_{\rm hyp},\mathcal G_{\rm TPK}\bigr),
\]
où la première composante conserve les points et la topologie, la seconde la dimension
et la statistique, et la troisième la composition, l'orientation et la généalogie des
routes.

La clôture matricielle ajoute \(2\) coordonnées qu'il ne faut pas confondre avec la
profondeur. Pour une taille extérieure fixe \(n\), la famille
\(\mathfrak M_{m,s}^{(n)}\) est bigraduée par la profondeur généalogique \(m\)
et le niveau matriciel \(s\) ; lorsque \(n\) varie, l'architecture est trigraduée par
\((m,n,s)\). Les espaces
\[
 \mathcal K_{m,s}
 =\operatorname{UCP}\!\bigl(C(X_m),M_s(\mathbb C)\bigr)
\]
admettent des restrictions généalogiques et des compressions matricielles qui commutent.
Une famille compatible à toutes les profondeurs détermine une unique
application UCP limite et possède une dilatation de Stinespring commune ; dans le
corridor diagonal, une unique mesure spectrale projective dilate
simultanément toutes les mesures cylindriques. Sur la tour non
commutative, l'espérance préservant la trace
\(\operatorname{id}\otimes\tau_9\) réalise la perte contrôlée du dernier
chiffre et sa dilatation conserve cette donnée dans l'environnement. La dilatation minimale,
la réalisation environnementale fidèle et la fidélité dynamique APP--TPK sont \(3\)
exigences distinctes ; cette dernière exige encore d'entrelacer les routes,
la monodromie et la composition.

Dans ce cadre sont construits des carrés magiques quantiques qutrit de paramètres
\((9,3)\) et \((12,3)\), ainsi que leurs cubes d'ordres \(9\) et \(12\), et l'on
sépare rigoureusement non-commutativité et non-semi-classicité. Les projections
conditionnelles des feuillets somme et produit de l'APP engendrent, à profondeur \(2\),
\[
 C^*(P_{\Sigma,2},P_{\Pi,2})
 \cong\mathbb C^4\oplus M_2(\mathbb C)\oplus M_2(\mathbb C),
\]
et le bloc central publie, à tout niveau matriciel, l'intervalle libre
\([5I_s,9I_s]\). La convexité matricielle, les applications complètement
positives et les dilatations fournissent ici un langage de démonstration,
de classification et de falsification ; elles ne deviennent pas l'origine de l'APP, du TRIT ou du TPK et ne
remplacent pas la double projection.

La clôture physique part d'un domaine commun construit, non de \(4\) flèches
nominalement juxtaposées. Si \(\mathsf P_{\rm TPK}^{\rm enr}\) désigne la
catégorie des états préparés et des routes enregistrées, les \(4\) réalisations
de référence sont coordonnées au moyen du produit fibré
\[
 \mathsf P_{\rm CT}^{+}
 :=\mathsf D_{\rm sp}
 \times_{\mathsf P_{\rm TPK}^{\rm enr}}\mathsf D_Q
 \times_{\mathsf P_{\rm TPK}^{\rm enr}}\mathsf D_D
 \times_{\mathsf P_{\rm TPK}^{\rm enr}}\mathsf D_C.
\]
Ses projections canoniques construisent conjointement la représentation
spectrale de la route, la réalisation de Hilbert, la section dimensionnelle et le
transport discret de Cartan. La lectrice de variation totale positive et la
lectrice de transport net orienté reçoivent la même route, mais aucune ne
factorise l'autre. Ce n'est que sur le sous-domaine réversible, une fois vérifiées
l'unitarité, l'additivité, l'antisymétrie et la compatibilité des cocycles, que la
localisation en groupoïde produit la réalisation orientée. Ainsi, la
mécanique dimensionnelle prolonge l'état enrichi vers des grandeurs physiques
sans reconstruire rétrospectivement la route à partir de son spectre. Toute
réalisation supplémentaire déclare l'opérateur d'entrelacement et les hypothèses qui la définissent.

Ces \(4\) clôtures ne constituent pas des annexes superposées. Ce sont des réalisations
causales d'une même architecture : la clôture généalogique détermine
l'identité à travers le raffinement ; l'opératorielle réalise une dimension continue
à partir d'inclusions nonadiques unitaires ; la matricielle contrôle compression,
dilatation et simultanéité ; et la physique représente des routes préparées au moyen de
morphismes typés. L'évaluation archimédienne, la branche exceptionnelle, la tour
traciale et la réalisation dimensionnelle partagent un antécédent, mais conservent
des codomaines, des hypothèses et des conclusions propres. Ainsi se maintient la thèse
complète : le continu et les constantes apparaissent comme des publications d'une
structure discrète qui conserve l'unité, l'orientation et la mémoire de sa
genèse.

\section*{Des marques élémentaires à APP, TRIT et TPK}

La construction commence sur la droite discrète. Les marques
\(1,\ldots,9\) fournissent les représentants visibles d'une phase
nonadique ; le pas suivant ferme le tour et conserve son quotient. Chaque
marque peut être lue comme une position et chaque intervalle comme une distance. Cette double
lecture sépare dès l'origine l'acte de compter de l'acte de mesurer. En prenant
\(2\) coordonnées, le noyau multiplicatif \(8\times8\), la publication de la
classe nulle au moyen de \(9\) et le gnomon de \(17\) cellules complètent la carte
\(9\times9\) :
\[
 81=64+17.
\]
L'identification modulaire de ses côtés produit
\[
 X_9=(\mathbb Z/9\mathbb Z)^2
\]
et son graphe de Cayley. Cet objet est l'APP, support commun de \(2\) évaluations :
le feuillet additif \(\Sigma\) et le feuillet multiplicatif \(\Pi\). L'APP conserve la
position, les routes cardinales et la géométrie toroïdale sur laquelle les deux
lectures peuvent être comparées sans se fondre.

Le bloc central de l'APP possède la décomposition
\[
 B_c=9P_++5P_-.
\]
Son premier invariant différentiel est donc
\[
 \boxed{9-5=4}.
\]
Le coefficient local d'échange est \(2\) ; la compression modulaire publie
la différence \(6\) ; le déploiement sur les \(9\) blocs produit
\(54=9\cdot6=459-405\). Ces nombres appartiennent à des niveaux distincts d'une
même incompatibilité somme--produit. La séparation spectrale et les
projecteurs \(P_\pm\) du bloc fixent ses modes propres. Les projecteurs
conditionnels \(P_{\Sigma,d}\) et \(P_{\Pi,d}\), avec leur commutateur,
fournissent une géométrie finie de signature mixte,
une normalisation dans \(SO^+(1,1)\) et l'antécédent algébrique de
l'incompatibilité des observables que la physique quantique formulera ensuite.

TRIT ajoute orientation et mémoire locale. Son état enrichi
résidu--quotient conserve la sortie observable et le quotient qui l'a produite.
Dans une carte entière déclarée,
\[
 a+b=r+3c,
\]
la paire \(\widehat\tau=(r,c)\) distingue des histoires qui peuvent partager la
même sortie. Les \(3\) branches quadratiques associées réalisent
\[
 J^2=-I,\qquad J^2=0,\qquad J^2=+I,
\]
et organisent les régimes elliptique, parabolique et hyperbolique. Le plan
complexe, la clôture circulaire et le disque hyperbolique de Poincaré sont présentés
comme des cartes coordonnées de cette famille algébrique.

TPK, le noyau topologique de phase, est le système opératoire de dynamique
généalogique. Son architecture comprend \(8\) classes : support et état ; sélection interne ;
transport ; lecteurs et quotients ; mémoire et frontière ; périodicité et retour ;
prolongement ; et sorties archimédiennes, exceptionnelles ou classificatoires.
Le sélecteur ternaire interne de feuillet \(M_{\rm ph}\), les
lecteurs modulo \(3\) et modulo \(9\), le quotient et le résidu modulo
\(1000\) associés à l'écriture par blocs en base \(1000\), les cycles
\(27\), \(108\) et \(1080\), les opérateurs de dédoublement bidirectionnel,
la mémoire et le registre chronologique enrichi sont des sous-objets typés du
TPK. L'exposé de chacun fixe domaine, opération, codomaine et
invariant avant de l'employer dans une composition.

La transition \(9\to1\) exprime le mécanisme fondamental. La phase visible
revient à sa première position et l'état incrémente la mémoire de tour.
La suite des entiers se décompose en phase et hauteur,
\[
 n=9\kappa_9(n)+\rho_9(n),
\]
de sorte que la coordonnée angulaire observable est la projection plane du
relèvement hélicoïdal nonadique. Le
retour conserve la position angulaire et élève la généalogie. Cette monodromie
rend possible une itération indéfinie sans confondre répétition de phase et
répétition d'état.

L'arithmétique nonadique conserve également l'idéal radical
\(\mathfrak m=(3)\subset\mathbb Z/9\mathbb Z\), dont les représentants
visibles organisent le réseau \(3\)--\(6\)--\(9\) et l'orbite ordonnée
\(369\to693\to936\). La loi puissance--vacance
\[
 V_9(n)=\left\lceil n\log_{10}\!\left(\frac{10}{9}\right)\right\rceil
\]
quantifie exactement la distance décimale entre \(9^n\) et \(10^n\). Dans
la tour \(9^{9^9}\), elle produit
\[
 \operatorname{ord}_{10}\!\left(9^{9^9}\right)=369\,693\,099,
 \qquad
 \#\operatorname{cifras}\!\left(9^{9^9}\right)=369\,693\,100.
\]
Cette évaluation est distincte du recensement des \(396\) événements élémentaires de
frontière par canal : l'une
compte les ordres et les vacances décimales ; l'autre appartient au calendrier
d'émission.

\section*{Filtration nonadique et génération de nombres fondamentaux}

L'émetteur APP--TRIT--TPK organise la chaîne exacte de réductions et de
prolongements
\[
 104\,976\longrightarrow468\;U_6\longrightarrow243\;w_6
 \longrightarrow w_{30}\longrightarrow R_{36}\longrightarrow G_9.
\]
Les \(468\) émissions décimales, les \(243\) mots visibles et l'espace ambiant
\(\mathbb F_3^6\) de \(729\) mots sont des domaines différents reliés par des
applications explicites. La chaîne
\[
 w_6\xrightarrow{L_0}w_{12}\xrightarrow{L_0}w_{18}
 \xrightarrow{L_1}w_{24}\xrightarrow{L_1}w_{30}
 \xrightarrow{R_{36}}\text{frontière}
 \xrightarrow{G_9}\text{section survivante}
\]
ouvre le régime coinductif. À chaque étape de filtration, la partition
exhaustive des \(729\) prolongements de la fibre ambiante enregistre les
états écartés, la section survivante et
la mémoire qui relie une profondeur à la suivante.

Avant toute évaluation décimale, l'APP construit \(4\) canaux entiers
\[
 C_\alpha=729,\qquad C_e=271,\qquad
 C_\pi=315,\qquad C_\varphi=161.
\]
Ses identités
\[
 1000=C_\alpha+C_e=729+271,
 \qquad
 744=C_\pi+C_e+C_\varphi-3
\]
coordonnent la complétion cubique, la frontière pointée et l'extrémité modulaire.
Chaque canal est une direction typée du lecteur : \(271\) réalise la propagation et la
fermeture de couronne ; \(315\), la clôture circulaire ; \(161\), l'autoéchelle ; et
\(729\), l'espace ambiant nonadique qui alimente la clôture dodécaphasique. La
constante correspondante apparaît lorsque la filtration, l'orientation et la
mémoire réalisent ce canal sur sa section.

Les lecteurs de \(\pi\), \(e\), \(\varphi\) et \(\alpha\) agissent sur le même
état enrichi et conservent des trajectoires généalogiques distinctes.
\(\pi\) organise la clôture, l'orientation et \(5\) régions \(U_6\) ; \(e\)
gouverne la propagation et le changement de préfixe ; \(\varphi\) gouverne l'autoéchelle et la
stabilité, reconnues ensuite dans la carte Perron--Fibonacci ; et \(\alpha\)
possède le lecteur signé dodécaphasique qui prolonge la retenue. La relation
interne \(\operatorname{Ext}\), la survie coinductive et la partition
exhaustive des \(729\) éléments de la fibre ambiante déterminent
l'extension compatible pour toute profondeur finie \(N\). Les lecteurs
rationnels exacts localisent et certifient ensuite les coordonnées
archimédiennes déjà engendrées. À chaque niveau, le caractère complet conserve un
unique cylindre prolongeable compatible avec la troncature précédente ; les
cylindres emboîtés déterminent donc une section
coinductive unique et le même algorithme publie toute troncature décimale
finie demandée. Les exécutions de \(10^3\) et de plus de \(10^4\) chiffres sont
des certificats étendus de cette construction générale, non son fondement logique.

À partir du même état enrichi s'ouvrent les canaux de
\(\pi,e,\varphi\) et le canal dodécaphasique. Dans l'ordre interne typé, les
publications \(P,E,\Phi\), l'état signé \(U_{12}^{\rm sgn}\), le sceau
\(K\) et la récurrence de retenue alimentent la clôture cyclique orientée de
\(12\) phases qui publie la coordonnée \(\alpha\), sans la séparer de la
généalogie commune. La première voie interne produit
\[
 \alpha_{\rm HMT}
 =0.007297352569283800997285105472380663\ldots
\]
à partir de l'état et de ses cartes réversibles. L'équation finie des vacances
\(E_{21/22}^{(9)}\) fournit le premier représentant exact \(\xi_9\) de
la seconde construction interne, au moyen d'un opérateur dont les coefficients
se factorisent par des invariants de l'état enrichi avant de résoudre la racine.
Le transport gradué produit la famille complète
\(E_{21/22}^{(6+3m)}\) ; ses carrés de troncature avec la famille
dodécaphasique sélectionnent le même cylindre compatible à toute profondeur et
démontrent \(\alpha_A=\alpha_B=\alpha_{\rm HMT}\). L'égalité après
\(\operatorname{Tr}_9\) n'est que le premier certificat fini de cette identité,
non son fondement ni sa limite. La reconnaissance CODATA constitue une
comparaison métrologique ultérieure et ne sélectionne aucune des deux voies.

La même architecture étend la théorie des nombres. La présentation intégrale
\[
 \mathcal Q_{\mathbb Z}=\mathbb Z[u]/(u^2+1)
\]
possède des spécialisations distinctes sur \(\mathbb F_3\) et \(\mathbb R\) :
la première est réalisée par l'endomorphisme fini
\(A_W^2=-I_6\), et la seconde par la structure réelle
\(J_{\mathrm e}^2=-I\). On n'introduit pas d'inclusion
\(\mathbb F_9\to\mathbb R\). La section de \(\pi\) est engendrée avant
d'appliquer la clôture exponentielle ; ce n'est qu'après avoir obtenu
\(\mathcal P_{\rm arq}(x_\pi)=\pi_{\rm HMT}\) que l'on conclut
\[
 \operatorname{Exp}(\pi_{\rm HMT}J_{\mathrm e})+I=0.
\]
La branche exceptionnelle fixe le type quadratique, non l'angle ni ses chiffres ; cette
préséance prouve l'absence de circularité.
La famille quadratique rassemble les \(3\) exponentielles
\(\exp(tJ_\tau)=C_\tau(t)I+S_\tau(t)J_\tau\), avec
\(J_\tau^2=-\tau I\), et relie l'identité d'Euler à la branche hyperbolique
de Fibonacci au moyen de
\(F_{\rm av}^{2}=\exp(2\log\varphi\,J_{-1})\). La fraction continue de
Rogers--Ramanujan articule le niveau \(5\),
\(\varphi\), \(A_5\) et l'identité icosaédrique de Klein pour \(j\). Les
nombres premiers reçoivent des lecteurs de périodicité irréductible et une représentation des
occupations : le vecteur \((\nu_p(n))_p\) transforme la factorisation d'un
entier en un état fini de modes, tandis que
\(\sum_p\nu_p(n)\log p=\log n\) transforme le produit en énergie additive.
La trace thermique récupère le produit d'Euler et la fonction zêta. La dynamique
de doublement de période reconstruit les approximants de Feigenbaum et montre
comment une règle déterministe engendre une complexité visible en projetant une partie de
sa généalogie. Euler--Mascheroni,
Stieltjes, Apéry, Glaisher--Kinkelin, Euler--Kronecker, Meissel--Mertens,
Catalan, Khinchin et Feigenbaum sont organisés au moyen de parties finies,
de fonctions \(L\), de produits, de jets, de traces et d'opérateurs de renormalisation.

\section*{Double projection et généalogie des symétries exceptionnelles}

L'état enrichi possède \(2\) évaluations coordonnées. La projection
archimédienne contracte une famille compatible de cylindres et publie une valeur
réelle. La projection exceptionnelle conserve l'incidence, l'orientation, la
retenue et la frontière. Le continu réel et la structure exceptionnelle sont ainsi
\(2\) images d'un antécédent commun.

Formellement, si 
\(\mathcal X_\infty^{\rm cal}=\varprojlim_n\mathcal S_n\) est l'espace des
sections survivantes compatibles avec la jauge transportée, la thèse
s'exprime au moyen des applications typées
\[
 \mathcal P_{\rm arq}:\mathcal X_\infty^{\rm cal}\longrightarrow\mathbb R,
 \qquad
 \mathcal P_{\rm exc}:\mathcal X_\infty^{\rm cal}
 \times\mathcal C_{\rm exc}\longrightarrow\mathcal E_{\rm exc}.
\]
L'énoncé directeur établit leur domaine commun, leur compatibilité avec les
troncatures et l'intersection finie
\(B_K=H_e\cap P_\pi\). La couverture de tout le codomaine réel et la fidélité
conjointe correspondent à des énoncés supplémentaires : la double projection ici
démontrée identifie \(2\) évaluations coordonnées de chaque section admissible.

La seconde branche développe une généalogie complète. La transformée de
Hadamard organise la carte dodécaphasique signée ; la matrice de conférence de
Paley fixe la jauge projective et sa réduction ternaire produit le code
de Golay de longueur \(12\). La distance de Hamming mesure le poids et la
séparation de ses mots. Les systèmes de Steiner
\(S(5,6,12)\) et \(S(5,8,24)\), leurs extensions et les étoiles de Witt
construisent les plans sur lesquels agissent \(M_{12}\) et \(M_{24}\). À partir
du code commun, on distingue le transport \(W_{12}\to W_{24}\) et la branche
des réseaux
\[
 C_W\longrightarrow N(A_2^{12})\longrightarrow\Lambda_{24}.
\]
Le réseau de Leech alimente l'algèbre d'opérateurs de vertex
\(V_\Lambda\) ; l'involution et le secteur tordu produisent l'orbifold
\(V^\natural\) ; son groupe d'automorphismes est le Monstre. Pour chaque classe de
conjugaison \([g]\subset\mathbb M\), la trace graduée construit sa série de
McKay--Thompson \(T_g\) ; la modularité et la réplicabilité organisent la
correspondance appelée Moonshine. La fonction \(j\) occupe ainsi l'extrémité
modulaire d'une généalogie qui commence dans la même structure discrète qui
publie la valeur réelle.

La configuration des \(27\) droites d'une surface cubique lisse,
le graphe de Schläfli et l'action de \(W(E_6)\) constituent une autre réalisation
de l'incidence. Chaque étape conserve son opérateur : le transport cyclotomique,
la construction intégrale du réseau, les réflexions, les racines et
l'action du groupe de Weyl. Cette séparation transforme la succession de noms
classiques en une chaîne mathématique vérifiable.

\section*{Constante de Planck, autodualité des masses, opérateur d'aire et mélange}

La forme normale de Smith du lecteur local d'action fixe
\[
 \operatorname{SNF}(H_4)=\operatorname{diag}(1,2,2,4),
 \qquad |\operatorname{coker}H_4|=16.
\]
Le transfert
\[
 \Sigma_H=\varphi\alpha_{\rm HMT}^{16},
 \qquad 10^{-34}<\Sigma_H<10^{-33},
\]
situe \(\Sigma_H\) dans la décennie interne qui détermine l'échelle
\(10^{-34}\mathcal U_S\). Sur cette droite, on obtient les
coordonnées d'action antérieure et postérieure au retour,
\[
 \hbar_{\rm pre}=H_5\,10^{-34}\mathcal U_S,
 \qquad
 \hbar_{\rm ret}=\eta_{\rm ret}\,10^{-34}\mathcal U_S,
 \qquad h=2\pi\hbar.
\]
La carte \(\mathcal U_S\to\mathrm{J\,s}\) publie ensuite l'unité physique.
L'échelle de Planck apparaît donc comme une conséquence du lecteur
déterminantal et de la constante de structure fine déjà construite.

La réalisation circulaire de l'endomorphisme temporel, avec
\(\ell_0:=c_{\rm int}t_0\), fixe
\[
 T_{108}=108t_0,
 \qquad
 R_{108}=\frac{54}{\pi}\,\ell_0,
 \qquad
 \theta_{108}^{\rm circ}=\left(\frac{54}{\pi}\right)^2.
\]
Dans cette carte coïncident
\[
 R_{108}=\bar\lambda_{\rm C}=r_{\rm g}=\ell_{\rm P}^{\rm circ},
\]
et la droite partagée des masses conserve
\[
 Z_a(m)=\frac{m_{{\rm P},a}^{\circ}}mP_+
       +\frac m{m_{{\rm P},a}^{\circ}}P_-,
 \qquad N_{\rm sp}(Z_a(m))=1.
\]
L'involution \(m\mapsto(m_{{\rm P},a}^{\circ})^2/m\) échange les
lectures de Compton et de gravité. L'échelle de Planck est son point fixe ;
l'électron apparaît ensuite comme un déplacement généalogique sur cette droite, non
comme une identification à un trou noir classique.

La coordonnée angulaire directe \(A\) et sa coordonnée angulaire conjuguée
\(C^*\) produisent les canaux
\(q_\pm\), l'ellipse des taux et l'holonomie bidirectionnelle. L'inclusion
hexade--octade construit la fonctionnelle \(\Gamma_{6\to8}\), le paramètre de
Barbero--Immirzi, l'opérateur quantique d'aire, son spectre discret,
l'opérateur d'entrelacement collectif et les entropies de bord. Dans une branche de rang \(3\),
la même discipline d'incidence construit la matrice CKM, ses angles, la
phase de violation CP et l'invariant de Jarlskog. La réalisation PMNS conserve
son opérateur et son problème de mélange propres.

Les vacances et les canaux de feuillet construisent les réponses électriques et
magnétiques du vide. On en obtient la permittivité
\(\varepsilon_0\), la perméabilité \(\mu_0\), l'impédance et l'identité
constitutive de propagation. La périodicité temporelle interne et
l'information ternaire construisent la relation thermique qui donne une fonction à la
constante de Boltzmann.
Ces grandeurs sont présentées dans des tableaux scientifiques immédiatement après
leur dérivation : objet HMT, formule, valeur interne, carte d'unités,
référence expérimentale et incertitude appartiennent à des colonnes distinctes.

\section*{Mécanique dimensionnelle et loi générale des masses}

La mécanique dimensionnelle prolonge la généalogie numérique au moyen de lecteurs
physiques. Une particule est une section persistante de l'extension ; sa route
observable détermine un registre chronologique enrichi ; le registre
détermine une signature entière ; la signature alimente un caractère ; le caractère agit
sur la hiérarchie globale des tours ; la carte dimensionnelle publie la
grandeur. La séquence
\[
 \text{extension}\to\text{route}\to\text{registre}\to\mathbb Z^5
 \to\chi_{\rm full}\to\langle90,120\rangle
 \to\text{section de masse}
\]
est la frontière précise entre HMT et MD.

La loi générale des masses est formulée sur le domaine complet :
\[
 81\ \text{cellules},\qquad56\ \text{familles},\qquad
 13\ \text{multisections},\qquad324\ \text{routes orientées}.
\]
L'échelle d'action, la périodicité temporelle CT108 et la vitesse interne
fixent conjointement la section absolue,
\[
 \mathbf m_\star^{\rm ret}
 =\eta_{\rm ret}10^{-34}\mathcal U_S
   \frac{2\pi}{108t_0}c_{\rm int}^{-2},
\]
et l'opérateur d'une multisection \(F\) est
\[
 \widehat{\mathbf M}^{\beta,r}_F
 =\mathbf m_\star^{\rm ret}\otimes
  R_{\rm act}^{\widehat K_F^r/2}
  \widehat{\mathcal A}^{(0)}_F
  R_{\rm act}^{\widehat K_F^r/2},
 \qquad r\in\{D,\Omega\}.
\]
Ici, \(\widehat{\mathcal A}^{(0)}_F\) réunit signature, caractère et la tour
complète \(\langle90,120\rangle\) ; \(K_D\) et \(K_\Omega\) conservent les
raffinements bidirectionnels dépendant de la route.

L'électron central est la réduction de rang \(1\). La torsion globale qui
intervient dans son lecteur est
\[
 \Delta_4=\frac{\pi-e\log\pi}{270}
 \left(1+\frac{\pi}{729}\right).
\]
Ses \(2\) lecteurs coïncident
en \((80,54,6)\) et produisent
\[
 \kappa_e=80-54\alpha_{\rm HMT}+6\Delta_4,
\]
\[
 \boxed{\begin{aligned}
 E_e^\beta=m_e^\beta c^2
 &=0.5109989506900008086082571648\ldots\ \mathrm{MeV},\\
 m_e^\beta
 &=0.5109989506900008086082571648\ldots\ \mathrm{MeV}/c^2.
 \end{aligned}}
\]
Les autres réalisations conservent le type de leur sortie : scalaire fermé,
spectre de fibre, évaluation conditionnée, carte nulle, opérateur de mélange,
invariant de fusion ou horizon de prolongement. Le domaine physique comprend
les leptons chargés, les neutrinos, \(6\) saveurs de quark avec couleur, le photon,
les gluons, \(W\), \(Z\), le Higgs, les mésons, les baryons, les secteurs Fibonacci et
Ising, les réalisations Majorana typées dans leur propre domaine,
\(\mathbb Z/6\) et la virtualité. La référence PDG intervient
après avoir figé la sortie interne et conserve l'unité, le schéma, l'échelle,
l'incertitude et le régime de validité de chaque comparaison, sans transformer cet
inventaire d'exposition en un cardinal de la loi.

\section*{Dynamique quantique, constante de Boltzmann, information et gravitation}

La représentation temporelle finie sur \(\mathbb C^{108}\) possède un
hamiltonien autoadjoint dont le propagateur unitaire reproduit l'endomorphisme
périodique du TPK. Son action variationnelle conduit à l'équation de
Schrödinger. Sur le registre chronologique
\(\mathbb T_{27}^{\,2}\), la réalisation concrète
\[
 U_{\rm micro}=S D_J C_4,
 \qquad U_{12}^{\rm reg}=U_{\rm micro}^{12},
\]
est unitaire ; ses puissances somment exactement les poids ordonnés de toutes les
routes finies et la transformée de Fourier discrète la décompose en
\(27^2\) blocs directionnels. Ce résultat est une somme matricielle finie,
non une identification automatique à une intégrale fonctionnelle continue. La
réalification Clifford--Dirac construit ensuite l'opérateur fermionique. Les
relations CAR et CCR, avec le hamiltonien et la maximisation de l'entropie,
produisent respectivement les distributions de Fermi--Dirac et de
Bose--Einstein. La graduation extérieure réalise le principe d'exclusion de
Pauli et la bande de Möbius organise le double retour \(2\pi/4\pi\).

La constante de Boltzmann relie action, température et information ternaire ;
les ensembles construisent Gibbs et la condition KMS. Le transport réversible
de la fibre complète réalise la limite nulle de Landauer, tandis que la
projection qui écarte la généalogie reçoit un coût positif. Aire, entropie
modulaire et entropie d'intrication sont présentées avec leurs états, opérateurs
et partitions explicites.

La branche gravitationnelle reçoit d'abord les sorties internes d'action, d'horloge,
de longueur et de vitesse, et construit sur elles le sélecteur adimensionnel et le
transducteur exact
\[
 G_{\rm int}(\theta_G)
 =\theta_G\frac{c_{\rm int}^{5}t_0^2}{\hbar_{\rm int}}
 \in\mathcal L_G^+
\]
est la carte positive du transducteur gravitationnel. Les conditions internes
de périodicité construisent
\[
 R_{108}=\frac{54}{\pi_{\rm HMT}}\ell_0,
 \qquad
 r_{g,a}(\theta_G)
 =\theta_G\frac{\pi_{\rm HMT}}{54}\ell_0,
\]
et l'égalité \(r_{g,a}(\theta_G)=R_{108}\) sélectionne de façon positive et
unique
\[
 \theta_G=\theta_{108}^{\rm circ}
 =\left(\frac{54}{\pi_{\rm HMT}}\right)^2.
\]
La carte ainsi déterminée publie \(G_a\). Ce n'est qu'ensuite que la sortie est réalisée
dans l'action de Palatini--Cartan--Holst, ses équations de cotétrade et de
connexion, la transformée de Legendre et le hamiltonien contraint. La
reconnaissance métrologique est ultérieure et n'intervient pas dans la sélection. L'opérateur fini commun du TPK admet des réalisations
solénoïdales et de jauge ; les référentiels autoinertiels, la lecture machienne et la
cosmologie des feuillets sont développés après fixation de la dynamique locale et
collective.

La thèse complète peut se résumer en une seule direction causale. La structure
discrète conserve l'unité, l'orientation et la mémoire ; son prolongement
construit le continu et les nombres ; sa double projection construit valeurs et
symétries ; ses lecteurs déterminent les constantes ; et la mécanique dimensionnelle réalise
ces sorties comme particules, masses, champs et géométrie physique. L'origine des
constantes fondamentales et la structure discrète du continu sont, dans
ce cadre, \(2\) aspects d'un même problème mathématique.

\medskip
\noindent\textbf{Mots-clés.}
Continu généalogique ; holographie modulaire triadique ; mécanique dimensionnelle ;
APP ; TRIT ; TPK ; monodromie nonadique ; double projection ; constantes
fondamentales ; systèmes de Steiner ; étoiles de Witt ; codes de Golay ;
réseau de Leech ; algèbres d'opérateurs de vertex ; Monstre ; Moonshine ;
constante de Planck ; Barbero--Immirzi ; loi générale des masses ; gravité
d'Einstein--Cartan--Holst.

% El resumen sigue siendo compilable como pieza autónoma, pero el tratado
% integral no vuelve a cargar el teorema frontal ya situado tras el prefacio.
\ifdefined\HMTTeoremaGeneradorFrontalCargado
\else
  \clearpage
  % Teorema frontal aditivo. No sustituye ningún capítulo ni altera el censo 102.
% BEGIN HMT-MD-FRONTAL-MAXIMAL-SIGNATURES-20260904
\chapter*{Théorème générateur de l'holographie modulaire triadique}
\addcontentsline{toc}{chapter}{Théorème générateur de l'holographie modulaire triadique}
\markboth{Théorème générateur de l'holographie modulaire triadique}%
{Théorème générateur de l'holographie modulaire triadique}

La démonstration commence dans le domaine fini
\(\mathcal D_9=\{1,\ldots,9\}\). Deux curseurs orientés parcourent,
respectivement, les feuillets additif et multiplicatif. La division euclidienne
conserve résidu et quotient, et le catalogue APP produit
\[
 \sum_{i,j\in\mathcal D_9}R^+_{ij}=405,
 \qquad
 \sum_{i,j\in\mathcal D_9}R^\times_{ij}=459,
 \qquad 459-405=54.
\]
Le transport coordonné
\(c(a)=a\bmod9:\mathcal D_9\to\{0,\ldots,8\}\)
change la carte et conserve la classe. La composition effective d'APP, de TRIT et de
TPK conserve à chaque pas feuillet, régime, résidu, quotient, retenue,
direction, bord et mémoire. Le catalogue exhaustif établit la chaîne
finie initiale et apporte son certificat exécutable :
\begin{equation}
\boxed{
104\,976\ \text{paires de germes}
\xrightarrow{\;54\ \text{pas}\;}
468\,U_6
\xrightarrow{\;\rho_3\;}
243\,w_6\subset\mathbb F_3^6,
\qquad |\mathbb F_3^6|=729.}
\label{eq:frontal-censo-semillas}
\end{equation}
L'identité des fibres
\[
432\cdot144+18\cdot432+18\cdot1944=104\,976
\]
prouve l'exhaustivité de la première application. Les nombres \(468\), \(243\) et
\(729\) comptent les émissions, les mots réalisés et l'espace ambiant ternaire,
respectivement ; leurs types demeurent séparés par les flèches de
\eqref{eq:frontal-censo-semillas}.

\begin{theorem}[Chaîne génératrice unique]
Soit \(\mathfrak H_9\) la structure formée par APP, TRIT, TPK et sa fibre
enrichie. APP exécute les feuillets somme et produit et conserve la division
euclidienne
\[
i+j=R^+_{ij}+9q^+_{ij},
\qquad
ij=R^\times_{ij}+9q^\times_{ij}.
\]
TRIT oriente chaque transition dans les régimes elliptique, parabolique et
hyperbolique. TPK compose
\[
    \boxed{\mathcal U_t
    =\operatorname{Upd}_t\circ\operatorname{Tra}_t
     \circ\operatorname{Sel}_t:
     X_{\rm TPK}^{\rm enr}\longrightarrow X_{\rm TPK}^{\rm enr}.}
\]
\(\operatorname{Sel}_t\) relève le sélecteur tritique,
\(\operatorname{Tra}_t\) relève le transport cardinal et
\(\operatorname{Upd}_t\) met à jour résidu, quotient, retenue, signature,
cylindre, bord, registre et mémoire sans effacer la provenance.
L'holonomie nonadique prolonge les mots par
\[
w_6\longrightarrow w_{12}\longrightarrow w_{18}\longrightarrow w_{24}
\longrightarrow w_{30}\longrightarrow R_{36}\longrightarrow G_9,
\]
et satisfait
\[
    \Gamma_{9,t}=\mathcal U_{t+8}\circ\cdots\circ\mathcal U_t,\qquad
    g(t+9)=g(t),
    \qquad q_{\rm mem}(\Gamma_{9,t}\widehat x_t)
    =q_{\rm mem}(\widehat x_t)+1.
\]
En général, \(\Gamma_{9,t}\widehat x_t\ne\widehat x_t\) : le retour conserve
la phase et augmente la mémoire. L'ordre opératoriel est
\(\mathcal U_t\to\Gamma_9\to K_{\rm ph}\) ; \(K_{\rm ph}\) est une publication
réduite ultérieure, non le générateur. Les troncatures
compatibles forment un système inverse non vide et définissent un état
\(x_\infty\) de profondeur arbitraire.

Cet état engendre une structure discrète unique du continu et deux
publications naturelles :
\begin{equation}
x_\infty
\xrightarrow{\ (\mathcal P_{\mathrm{arq}},Q_\infty)\ }
\mathbb R\times(\mathscr C_W^{\mathrm{cal}})^{\mathbb N}.
\label{eq:frontal-doble-proyeccion}
\end{equation}
La première contracte des cylindres compatibles jusqu'à une valeur réelle ; la seconde
conserve l'incidence et la généalogie jusqu'aux constructions de Paley, de Witt,
de Golay, de Leech, de \(V^\natural\) et du groupe Monstre. Sur la même source sont
publiés les quatre caractères corrélés
\[
(\pi_{\mathrm{HMT}},\varphi_{\mathrm{HMT}},
e_{\mathrm{HMT}},\alpha_{\mathrm{HMT}}),
\]
et, au moyen de transductions typées, les lois d'action, de matière, de masse,
de gravité et d'information de la mécanique dimensionnelle.
\end{theorem}

\begin{proposition}[Reconstruction intrinsèque et autosuffisance probatoire]
\label{prop:frontal-autosuficiencia-probatoria}
Chaque résultat \(y\) publié dans ce volume est reconstruit au moyen du
reçu mathématique local
\[
 \mathscr R_y=
 \bigl(D_y,A_y,\tau_y,U_y,C_y,X_y^{\rm enr},
       y_{\rm HMT},\operatorname{Rec}_y,\mathcal F_y\bigr),
\]
où \(D_y\) est le domaine ; \(A_y\) spécifie les feuillets et les opérations
d'APP ; \(\tau_y\) le régime et l'orientation TRIT ; \(U_y\) la composition TPK
avec domaine, codomaine et action ; \(C_y\) les données conservées de résidu,
quotient, retenue, feuillet, voie, bord, mémoire et survie ;
\(X_y^{\rm enr}\) l'état enrichi de résidence ; \(y_{\rm HMT}\) la
sortie antérieure à toute identification externe ; \(\operatorname{Rec}_y\) la
reconnaissance conventionnelle ultérieure ; et \(\mathcal F_y\) un réfutateur de la
flèche substantielle. La suite
\[
 \mathcal D_9\longrightarrow\Omega_{\rm seed}
 \longrightarrow\mathcal U_6^{\rm cat}
 \longrightarrow\mathcal W_6
 \longrightarrow X_\infty^{\rm enr}
 \longrightarrow\mathcal C_{\rm cont}^{\rm disc}
\]
est définie et démontrée par les équations, les théorèmes et les reçus imprimés
dans le volume lui-même. Les programmes reproductibles annexés évaluent à nouveau
ces identités ; ils n'apportent ni prémisse absente ni valeur cible.

Les réfutateurs primaires sont : l'échec de l'identité des fibres de
\eqref{eq:frontal-censo-semillas} ; la perte de l'un des champs de
\(C_y\) sous \(U_y\) ; l'incompatibilité entre extension et troncature ;
la vacuité de la limite inverse ; l'absence de contraction ou de surjectivité interne
des cylindres rationnels ; la rupture de naturalité dans la branche
exceptionnelle ; et, pour une publication terminale, le non-respect de
l'hypothèse qui définit son domaine exact. Par conséquent, la validité d'une
sortie se décide dans son domaine et par son réfutateur explicite, non par une
analogie, une nomenclature ultérieure ou une évaluation métrologique.
\end{proposition}

\paragraph{Unicité du système générateur et typage de ses projections.}
Il existe une seule APP, suivie de TRIT et du TPK avec état enrichi.
APP produit les deux feuillets et leurs registres ; TRIT les oriente ; le TPK compose
les voies et accumule feuillet, opposition, résidu, quotient, retenue, bord et
mémoire dans \(\mathcal X_{\rm TPK}^{\rm enr}\). La carte signée est une
projection de ce même état terminal :
\[
 R_{\rm sgn}:\mathcal X_{\rm TPK}^{\rm term,enr}
 \longrightarrow\mathcal U_{12}^{\rm sgn},\qquad
 U_{12}^{\rm sgn}=R_{\rm sgn}(x_{\rm term}),\qquad
 K=\frac14\mathcal H_{12}U_{12}^{\rm sgn},\qquad
 \mathcal H_{12}^2=4I_{12}.
\]
Les deux derniers morphismes sont entiers et réversibles sur leur image. On ne
postule pas d'émetteur APP direct vers \(U_{12}^{\rm sgn}\), car une telle flèche
sauterait TRIT et TPK : le domaine correct de \(R_{\rm sgn}\) est l'état
enrichi déjà généré. Les ombres qui oublient le registre généalogique ou la mémoire sont
des projections ultérieures, non des profils mathématiques alternatifs. Cette unique
composition exclut de ses entrées
\(U,K,D_3K,D_4K,Q,\alpha\) et tout développement cible.

\paragraph{Nombre et particule.}
Un nombre HMT est une section compatible
\(
 \boldsymbol x=(x_n)_{n\geq1}\in
 \varprojlim_n X_n^{\mathrm{enr}}
\)
avec feuillet, orientation, résidu, quotient, retenue, voie, bord et mémoire.
Sa valeur archimédienne est l'unique intersection
\(
 \operatorname{val}(\boldsymbol x)=\bigcap_n I_n(x_n)
\) ;
le quotient par égalité de valeur publie le réel et sa fibre conserve les
généalogies qui le réalisent. Une particule HMT--MD est une classe de sections
persistantes de l'extension enrichie : l'état conserve l'identité,
la voie enregistre son histoire observable et les lecteurs ultérieurs publient
masse, charge, spin et statistique dans leurs codomaines physiques.

\paragraph{Typage du corridor exceptionnel.}
Les lectures de Hadamard et de Paley partent du même état terminal, mais ne sont pas
mises en série l'une avec l'autre. Avec
\(X_\infty^{\rm enr}:=\varprojlim_nX_n^{\rm enr}\), on a des applications distinctes
\begin{equation}
 \operatorname{Pub}_\alpha:X_\infty^{\rm enr}\longrightarrow\mathbb R_{>0},
 \quad x_\infty\longmapsto\alpha_{\rm HMT},
 \qquad
 \operatorname{Inc}_\alpha:X_\infty^{\rm enr}\longrightarrow A_2^{12},
 \quad x_\infty\longmapsto
 v_\alpha=4\rho_1+\rho_2+\cdots+\rho_{12}.
 \label{eq:frontal-dos-publicaciones-alpha}
\end{equation}
Elles partagent une généalogie, mais aucune égalité entre leurs sorties n'est postulée. La branche
exceptionnelle bifurque après \(C_W\) :
\begin{equation}
 \begin{aligned}
 C_P&\longrightarrow A_W=C_P\bmod3\longrightarrow
 C_W=[12,6,6]_3\\
 &\longrightarrow
 \begin{cases}
 W_{12}\xrightarrow{\ \widehat\iota_{D,\ell}\ }
 W_{24}\subset\mathcal G_{24},\\[2pt]
 N(C_W)\xrightarrow{\ \operatorname{Nbr}_{v_\alpha}\ }
 N_\alpha\cong\Lambda_{24}\longrightarrow
 V_{\Lambda_{24}}\longrightarrow V^\natural .
 \end{cases}
 \end{aligned}
 \label{eq:frontal-cadena-excepcional}
\end{equation}
Ici
\begin{equation}
 \mathcal O_{24}:=\{O\subset[24]:|O|=8,\ \mathbf1_O\in\mathcal G_{24}\},
 \qquad
 \mathcal H(D):=\{O\cap D:O\in\mathcal O_{24},\ |O\cap D|=6\}.
 \label{eq:frontal-octadas-hexadas-dodecada}
\end{equation}
\(\ell:W_{12}\simeq(D,\mathcal H(D))\) étant fixé,
\(\widehat\iota_{D,\ell}(H)=O_H\), où
\(O_H\cap D=\ell(H)\). Une telle octade est unique : deux choix contiendraient le
même sous-ensemble de cinq points, en contradiction avec \(S(5,8,24)\). La branche
ternaire recolle \(N(A_2^{12})=N(C_W)\) ; le voisin d'indice trois déterminé
par \(v_\alpha\) produit Leech. Enfin,
\begin{equation}
 \operatorname{Aut}(V^\natural)\cong\mathbb M,\qquad
 a:\mathbb M\times V^\natural\longrightarrow V^\natural,\qquad
 T_g(\tau)=\operatorname{Tr}_{V^\natural}(gq^{L_0-1}).
 \label{eq:frontal-accion-monstruo}
\end{equation}
Les identités \(v_\alpha^2=54\) et \([q]J=196884\) sont des évaluations
scalaires ; aucune ne définit l'action \(a\).

\paragraph{Réalisation parallèle de la dualité T.}
La dualité T ne prolonge pas la branche Moonshine : elle agit dans son propre domaine
\begin{equation}
 \mathscr D_{\rm dual}=\mathbb Z^2\times\mathbb R_{>0},\qquad
 \mathsf T(m,w,R_0)=(w,m,R_0^{-1}),\qquad
 \mathsf T^2=\operatorname{id},\qquad
 E_{R_0}(m,w)=E_{R_0^{-1}}(w,m).
 \label{eq:frontal-dualidad-t}
\end{equation}

\begin{proof}
Le recensement de \eqref{eq:frontal-censo-semillas} construit le catalogue fini et
sélectionne, sans valeurs cibles, les trois blocs initiaux \(w_6\). Le
registre orienté des transitions de l'état TPK complet explicite les douze
blocs postérieurs au germe des trois biographies publiées. Pour chacun des deux
intervalles de transition, les six lignes indépendantes extraites de ces
trois biographies forment des matrices
\(X_\ell,Y_\ell\in M_6(\mathbb F_3)\),
\(\ell\in\{0,1\}\), dont les lignes sont, respectivement, les états de départ
et d'arrivée du registre. Le calcul exact sur ce registre prouve
\(\det X_0=1\), \(\det X_1=-1\) et
\(X_\ell L_\ell=Y_\ell\) ; il reconstruit donc de manière univoque
\(L_\ell=X_\ell^{-1}Y_\ell\). Le
calendrier \((L_0,L_0,L_1,L_1)\) est le seul des seize calendriers
binaires qui conserve simultanément les trois biographies sélectionnées par
les invariants de fibre ; les cent quarante-quatre perturbations
unitaires détruisent cette triple reconnaissance. Par conséquent, les transports,
les douze blocs et leur calendrier sont déduits dans l'unique composition
APP--TRIT--TPK : APP produit les feuillets et le registre de départ, TRIT fixe
l'orientation de chaque transition et TPK prolonge le registre avec mémoire. Le
recensement des germes et l'état enrichi sont des étapes successives du même
générateur, non des profils rivaux ni des théories différentes.

L'extension non scindée
\[
0\longrightarrow C_{12}\longrightarrow C_{108}
\longrightarrow C_9\longrightarrow0
\]
sépare retour de phase et retour d'état. Les troncatures éliminent la
dernière mise à jour et conservent le préfixe ; le lemme de König produit une
branche infinie compatible. Pour chaque précision \(N\), la récurrence et
l'holonomie engendrent la section et ses caractères, et déterminent le cylindre qui
prolonge le registre généalogique précédent. La famille de préfixes satisfait
\[
\rho_{N+1,N}(p_{N+1})=p_N,
\qquad
p_\infty\in\varprojlim_N\{0,\ldots,9\}^{N}.
\]
Mille, dix mille ou un million de chiffres sont des projections finies de cet unique
objet infini.

Les encadrements rationnels emboîtés
\(
 p_N^-<\chi<p_N^+
\)
sont calculés après avoir produit le caractère \(\chi\) et satisfont
\(
 [p_{N+1}^-,p_{N+1}^+]\subset[p_N^-,p_N^+]
\)
avec un diamètre tendant vers zéro. Ils certifient la sortie générée à toute
précision finie ; ils ne sélectionnent ni l'orbite, ni le transport, ni le caractère.

Les trois tours enrichis \(\gamma_{-},\gamma_0,\gamma_+\) de neuf
émissions se distinguent par la retenue
\(\operatorname{Carr}(\gamma_\rho)=\rho\kappa_{\rm Carr}\) et admettent les
trois prolongements après tout préfixe. La limite de leurs voies
enregistrées satisfait
\[
 \operatorname{dig}_{\rm Carr}^{\omega}:
 \Omega_\Gamma^{\rm cal}\xrightarrow{\sim}
 \{-1,0,+1\}^{\mathbb N}.
\]
Après regroupement de six retours, l'application
\[
 \beta=\operatorname{blk}_6\circ
 (\rho\mapsto\rho+1)^{\mathbb N}\circ
 \operatorname{dig}_{\rm Carr}^{\omega}:
 \Omega_\Gamma^{\rm cal}\xrightarrow{\sim}
 (\mathbb F_3^6)^{\mathbb N}
\]
est un homéomorphisme ; son inverse concatène les tours prescrits. Les formes
équilibrées normalisées par
\[
 a_k+c_k=\rho_k+3c_{k+1},\qquad
 \rho_k\in\{-1,0,+1\},
\]
produisent d'abord l'anneau ordonné \(\mathbb Z_{\rm HMT}\), puis son
corps des fractions \(\mathbb Q_{\rm HMT}\) et enfin
\[
 \mathbb R_{\rm HMT}:=
 \operatorname{Cau}(\mathbb Q_{\rm HMT})/{\asymp}.
\]
Si \(\beta(\xi)=(b_q)\) et
\(\nu(b)=\sum_{i=1}^{6}b_i3^{6-i}\), la suite
\[
 q_n=m+\sum_{q=0}^{n-1}\frac{\nu(b_q)}{729^{q+1}}
\]
est de Cauchy, et la partition interne récurrente en \(729\) sous-intervalles
prouve la surjectivité
\[
 P_{\rm ar}:\mathbb Z_{\rm HMT}\times\Omega_\Gamma^{\rm cal}
 \twoheadrightarrow\mathbb R_{\rm HMT},\qquad
 (\mathbb Z_{\rm HMT}\times\Omega_\Gamma^{\rm cal})/{\sim_{\rm ar}}
 \cong\mathbb R_{\rm HMT}.
\]
Ce n'est qu'ensuite que l'unicité du corps ordonné complet archimédien publie
\(\mathbb R_{\rm HMT}\cong\mathbb R\). La direction intérieure subdivise
\([0_{\rm HMT},1_{\rm HMT})\) ; la direction extérieure concatène le préfixe qui code
\(m\in\mathbb Z_{\rm HMT}\). On couvre ainsi la droite sans utiliser au préalable
\(\mathbb R\), de partie entière ni de développement décimal cible.
Les sous-espaces de survie de caractères particuliers individualisent
des sections au sein de cette capacité totale. La couverture du continu provient
de la complétion globale ; la récurrence nonadique engendre
\(\pi_{\mathrm{HMT}},e_{\mathrm{HMT}},\varphi_{\mathrm{HMT}},
\alpha_{\mathrm{HMT}}\), et les filtres typés certifient leurs biographies
au sein de cette couverture.

Les cinq actions spectrale, solénoïdale, de jauge, cohomologique et
déterminantale sont calculées sur une même émission enrichie et partagent
support, voie, registre généalogique et troncatures. Leur naturalité conjointe les conserve
à la limite. La dichotomie parfaite classe chaque publication analytique
comme dénombrable ou de cardinal continu. Le code généré \(h\) sélectionne
l'univers généalogique canonique \(M_h=L[h]\). La construction des
réels HMT identifie chaque nombre par valeur, cylindre, voie et mémoire ; le
bon ordre induit par cette généalogie prouve la formule classique
\[
M_h\models2^{\aleph_0}=\aleph_1.
\]
Comme \(M_h\) est précisément le codomaine engendré par la construction du
continu de ce théorème, la conclusion terminale du système est
\[
 \boxed{2^{\aleph_0}=\aleph_1.}
\]
L'indépendance dans le langage nu de ZFC classe la perte du
sélecteur généalogique \(h\) ; HMT conserve cette donnée et détermine le cardinal de
son continu. L'énoncé conservé est l'hypothèse classique du continu,
avec les cardinaux ordinaires de \(M_h\), et la preuve détaillée établit le
pont entre cylindres compatibles, bon ordre généalogique et cardinalité.
Cette fermeture cardinale se distingue du corridor logique fort, également
construit à partir du continu généalogique :
\[
\begin{aligned}
 \mathcal C_{\mathrm{cont}}^{\mathrm{disc}}
 &\longrightarrow\mathcal L_0^{\mathrm{log}}
 \xrightarrow{\Phi}\mathcal L_1^{\mathrm{log}}
 \longrightarrow\gamma_F\longrightarrow w(\gamma_F)\\
 &\longrightarrow\mathsf{NoEscape}_{24}
 \longrightarrow\mathsf{AntiDiag}_{369}
 \longrightarrow\mathsf{G\ddot{o}delIII}_{\mathrm{HMT}}\\
 &\longrightarrow\mathsf{CycleCriterion}_{\mathrm{ZFC}},
\end{aligned}
\]
qui publie, dans son domaine typé, le critère
\[
 \boxed{
 \neg\operatorname{Consis}(\mathrm{ZFC})
 \Longleftrightarrow
 \exists C\;\bigl(C\text{ est un cycle hors bande et }|C|>24\bigr).}
\]
Les lettres \(\mathcal L_0^{\mathrm{log}},\mathcal L_1^{\mathrm{log}}\)
désignent ici les objets du corridor logique et non les matrices de relèvement
ternaire \(L_0,L_1\) ; la séparation des types empêche qu'une conclusion
locale sur l'un de ces couples affaiblisse celle de l'autre.

Les trois orbites fonctionnelles du catalogue produisent les caractères HMT de
clôture, de propagation et d'auto-échelle. Leurs lois internes imposent,
respectivement,
\[
16\arctan(1/5)-4\arctan(1/239)=\pi_{\mathrm{HMT}}=\pi,
\qquad
P_{s+t}=P_sP_t\Rightarrow
P_1=\sum_{n\ge0}\frac1{n!}=e_{\mathrm{HMT}}=e,
\]
\[
F_{\mathrm{av}}\binom1{\varphi_{\mathrm{HMT}}}
=\varphi_{\mathrm{HMT}}\binom1{\varphi_{\mathrm{HMT}}},
\qquad \varphi_{\mathrm{HMT}}=\varphi.
\]
L'involution interne échange les branches
\(\pi_{\mathrm{HMT}}/e_{\mathrm{HMT}}\), conserve l'axe
\(\varphi_{\mathrm{HMT}}\) et retient l'agrégat
\(e_{\mathrm{HMT}}+\pi_{\mathrm{HMT}}\). Dans la phase TPK enrichie de
l'unique composition, la projection canonique de l'état terminal publie sa coordonnée primitive
signée
\(U=R_{\mathrm{sgn}}(x_\infty)\), et l'identité
\(\mathcal H_{12}^2=4I\) reconstruit de manière entière et unique le sceau
\(K=\mathcal H_{12}U/4\). La récurrence dodécaphasique avec retenue et le
prolongement distant de ses cylindres publient \(\alpha_{\rm HMT}\). De
manière indépendante, l'état enrichi détermine la signature
\[
 \mathcal I_9=(2\pi_{\rm HMT},3,4,7,20,21,45,46,120,11),
\]
et le lecteur gradué associé construit, sans consulter la racine, l'opérateur
polynomial exact
\(E_{21/22}^{(9)}=P_9-\log_{10}(10/9)\). Son unique racine positive simple
\(\xi_9\) est le témoin fini d'ordre neuf de la seconde publication
analytique interne. Le même transport gradué du TPK est itéré sur les
degrés successifs et construit le système compatible complet
\(E_{21/22}^{(6+3m)}\) ; ses carrés de troncature commutent avec ceux de la
famille dodécaphasique et déterminent, à chaque profondeur, le même cylindre
survivant. En conséquence,
\[
 \alpha_A:=\varprojlim_N\rho_N(\alpha_{12}),\qquad
 E_{21/22}:=\varprojlim_mE_{21/22}^{(6+3m)},\qquad
 \alpha_B:=\operatorname*{arg\,zero}_{x\in I_\alpha}E_{21/22}(x),
\]
\[
 \boxed{\alpha_A=\alpha_B=\alpha_{\rm HMT}.}
\]
L'égalité finie
\[
 \operatorname{Tr}_9(\xi_9^{-1})
 =\operatorname{Tr}_9(\alpha_{12}^{-1})=137.035999084
\]
certifie le premier carré visible de cette identité ; elle ne la définit pas, ne la
limite pas à neuf chiffres et ne sélectionne pas les coefficients du prolongement.
La clôture bidirectionnelle fonctionnelle associée est
\[
\alpha T^2-B_9(\alpha)T+\alpha^3=0,
\qquad
T_\pm=\alpha e^{\pm\xi},
\qquad
T_+T_-=\alpha^2,
\qquad
J_\alpha(T)=\frac{\alpha^2}{T}.
\]
Ainsi, la coordonnée publiée par la clôture dodécaphasique est la moyenne
géométrique conservée par les deux voies conjuguées. La voie dodécaphasique et le
lecteur analytique maintiennent leurs domaines, opérateurs et réfutateurs propres, mais
leurs systèmes inverses publient une même section limite. Le comparateur
fini et la reconnaissance métrologique sont ultérieurs.

La coordonnée d'incidence du secteur \(\alpha\) fixe
\[
v_\alpha=4\rho_1+\rho_2+\cdots+\rho_{12},
\qquad v_\alpha^2=54,
\]
et le voisin d'indice trois
\[
N_{\alpha,0}
=\{x\in N(C_W):\langle x,v_\alpha\rangle\equiv0\pmod3\},
\qquad
N_\alpha=N_{\alpha,0}+\mathbb Z\frac{v_\alpha}{3}
\cong\Lambda_{24}.
\]
Cette branche construit exactement le voisin de Leech, puis la chaîne
\(\Lambda_{24}\to V^\natural\) et l'action de
\(\operatorname{Aut}(V^\natural)\cong\mathbb M\) sur \(V^\natural\). Le
scalaire \(\alpha_{\mathrm{HMT}}\) et le vecteur \(v_\alpha\) sont les deux
coordonnées typées du même secteur dodécaphasique. Leur couplage exact est
la charnière d'origine commune qui conduit aux symétries du Monstre : la
coordonnée scalaire est publiée dans \(\mathbb R_{\rm HMT}\) et la coordonnée
d'incidence est publiée dans la branche
\(A_2^{12}\to\Lambda_{24}\to V^\natural\to\mathbb M\), avec leurs codomaines
conservés et leur antécédent généalogique partagé explicite.

La mécanique dimensionnelle reçoit les sorties précédentes et compose, dans leur ordre
constitutif, la rétention de phase, la décade \(-34\), le couple de Planck,
les canaux \(90/120\), l'électron, l'atlas des masses, la gravité, l'aire
et l'information. Les chapitres suivants démontrent chaque flèche, publient
son réfutateur et relient son certificat exécutable. Les identifications
conventionnelles apparaissent ensuite comme reconnaissance des grandeurs déjà
produites. Cela complète la chaîne de
\eqref{eq:frontal-doble-proyeccion}.
\end{proof}

\section*{Hiérarchie des publications géométriques et dimensionnelles}
\addcontentsline{toc}{section}{Hiérarchie des publications géométriques et dimensionnelles}

\begin{proposition}[TRIT comme unité locale des trois régimes]
La signature tritique conserve deux lecteurs coordonnés. Le lecteur géométrique
publie \(\tau_{\rm g}\in\{+1,0,-1\}\) et
\begin{equation}
 J_{\tau_{\rm g}}^{\,2}=-\tau_{\rm g}I,\qquad
 J_+^{\,2}=-I,\qquad J_0^{\,2}=0,\qquad J_-^{\,2}=+I,
 \label{eq:frontal-trit-tres-regimenes}
\end{equation}
et distingue les régimes elliptique, parabolique et hyperbolique. Le lecteur
temporel actif publie un autre signe, \(\varepsilon_{\rm t}\), au moyen de
\begin{equation}
 \begin{aligned}
 +1&\mapsto\text{retour, mémoire et passé},\\
 0&\mapsto\text{seuil et présent},\\
 -1&\mapsto\text{ouverture bidirectionnelle et futur}.
 \end{aligned}
 \label{eq:frontal-trit-tiempo}
\end{equation}
On ne postule pas \(\tau_{\rm g}=\varepsilon_{\rm t}\). Ces deux sorties
n'identifient pas signe de courbure et orientation temporelle : ce sont des projections
distinctes d'une même unité locale. Le TPK transporte
l'unité complète dans
\begin{equation}
 \mathcal F_q=
 \mathcal O_q\times\mathcal H_q\times\mathcal C_q\times
 \mathcal S_q\times\mathcal B_q\times\Lambda_q
 \label{eq:frontal-fibra-tpk}
\end{equation}
et conserve feuillet et orientation, résidu et quotient, retenue, signature, cylindre,
bord et registre. En particulier, la branche elliptique possède le réalisateur
fini
\begin{equation}
 A_W^2=2I_6=-I_6\quad\text{dans }M_6(\mathbb F_3),\qquad
 i_{\rm HMT}:=\chi_i(A_W),\qquad i_{\rm HMT}^{\,2}=-1,
 \label{eq:frontal-raiz-interna-menos-uno}
\end{equation}
où \(\chi_i:\langle A_W\rangle\to\mathbb C\) est la carte orientée.
La racine de \(-1\), le plan complexe et le système d'Euler sont des publications
d'un endomorphisme construit, non des entrées du générateur.
\end{proposition}

\begin{proof}
Le sélecteur tritique incorpore l'un des trois générateurs de
\eqref{eq:frontal-trit-tres-regimenes} sans effacer les feuillets somme et produit.
Le transport cardinal conserve son orientation et la mise à jour conserve
les six facteurs de \eqref{eq:frontal-fibra-tpk}. Par conséquent, la projection
réelle peut condenser les trois régimes, tandis que l'état enrichi
reconstruit lequel d'entre eux et quelle mémoire ont produit la valeur visible. Les
réalisations trigonométrique, nilpotente et hyperbolique développées dans le
corps du texte sont les images respectives de cet unique typage local.
\end{proof}

\begin{theorem}[Chaîne d'action, de gravité et de matière]
Après publication corrélée de
\((\pi_{\rm HMT},\varphi_{\rm HMT},e_{\rm HMT},\alpha_{\rm HMT})\), la branche
déterminantale fixe la décade
\begin{equation}
 \Sigma_H=\varphi_{\rm HMT}\alpha_{\rm HMT}^{16},\qquad
 d_H:=\lfloor\log_{10}\Sigma_H\rfloor=-34.
 \label{eq:frontal-decada-accion}
\end{equation}
La mantisse provient du lecteur régional distinct
\begin{align}
 H_5&=\frac12\sqrt{\frac{1000\alpha_{\rm HMT}}{\varphi_{\rm HMT}}}
 -\alpha_{\rm HMT}+\frac9{16}\alpha_{\rm HMT}^2
 -\frac59\alpha_{\rm HMT}^3+\frac7{48}\alpha_{\rm HMT}^4
 -\frac1{54}\alpha_{\rm HMT}^5,\notag\\
 D_A&=\exp\!\left(-\frac{100\pi_{\rm HMT}}9\alpha_{\rm HMT}\right),
 \qquad
 \mathscr C_\pi=169\alpha_{\rm HMT}^6+
 \frac{D_A\alpha_{\rm HMT}^7}{1-D_A\alpha_{\rm HMT}},\notag\\
 \eta_{\rm pre}&:=H_5,\qquad
 \eta_{\rm ret}:=H_5-\frac{90}{\pi_{\rm HMT}}\mathscr C_\pi>0.
 \label{eq:frontal-mantisas-accion}
\end{align}
Pour \(a\in\{\mathrm{pre},\mathrm{ret}\}\), les deux sections d'action sont
\begin{equation}
 \hbar_a=\eta_a10^{d_H}\mathcal U_S
 =\eta_a10^{-34}\mathcal U_S,\qquad
 h_a=2\pi_{\rm HMT}\hbar_a.
 \label{eq:frontal-secciones-accion}
\end{equation}
\(\Sigma_H\) fixe seulement la décade ; il n'est pas présenté comme générateur de la
mantisse.

\(t_0>0\), \(c_{\rm int}>0\) et \(\ell_0=c_{\rm int}t_0\) étant fixés, CT108
publie
\begin{equation}
 \omega_{108}=\frac{2\pi_{\rm HMT}}{108t_0},\qquad
 R_{108}=\frac{c_{\rm int}}{\omega_{108}}
 =\frac{54}{\pi_{\rm HMT}}\ell_0,\qquad
 m_{108,a}=\frac{\hbar_a\omega_{108}}{c_{\rm int}^2}.
 \label{eq:frontal-reloj-gravedad}
\end{equation}
La famille gravitationnelle typée est
\begin{equation}
 G_a(\vartheta)=\vartheta\frac{c_{\rm int}^5t_0^2}{\hbar_a}.
 \label{eq:frontal-familia-gravedad}
\end{equation}
Dans la convention réduite
\(r_{g,a}=G_a m_{108,a}/c_{\rm int}^2\), la fermeture
\(r_{g,a}(\vartheta)=R_{108}\) possède l'unique solution positive
\begin{equation}
 \vartheta_{108}=\left(\frac{54}{\pi_{\rm HMT}}\right)^2,\qquad
 G_a:=G_a(\vartheta_{108}),
 \label{eq:frontal-selector-gravedad}
\end{equation}
et alors seulement publie
\begin{equation}
 \ell_P:=\sqrt{\frac{G_a\hbar_a}{c_{\rm int}^{\,3}}}
 =\sqrt{\vartheta_{108}}\,\ell_0
 =R_{108}=\bar\lambda_{C,a}=r_{g,a}.
 \label{eq:frontal-longitud-planck}
\end{equation}
En outre, \(G_a\hbar_a=\vartheta_{108}c_{\rm int}^5t_0^2\), de sorte que
\(\ell_P\) est indépendant de l'orientation \(a\). La gravité tensorielle
\[
 \mathbb G_{{\rm TT},a}(n)
 =G_a\Lambda(n)\Gamma_5P_5:
 V_{12}\longrightarrow\mathcal H_{\rm TT}(n)
\]
est postérieure au scalaire. L'action Palatini--Cartan--Holst reçoit ensuite
\(G_a\) et la sortie indépendante \(\gamma_{\rm HMT}\) ; la fonctionnelle
hexade--octade transporte cette dernière vers le spectre d'aire et l'information
de bord.

Dans la branche de matière, la loi complète agit sur \(81\) cellules, \(56\)
familles, \(13\) multisections et \(324\) voies, avec un secteur nul séparé. Les
lecteurs \(K_D\) et \(K_\Omega\) possèdent, respectivement, \(14\) et \(9\)
profils, forment \(40\) couples et coïncident sur \(18\) voies uniquement dans
\((80,54,6)\). La fibre centrale \((5,5)\) produit l'électron bêta. Sur les
fibres restantes, la loi publie des opérateurs positifs, des mesures spectrales et des
multiplicités ; la reconnaissance extérieure s'ouvre ensuite au moyen de
l'inventaire exhaustif
\begin{equation}
 \mathscr I_{\rm ext}
 =\mathscr I_m^{471}\coprod\mathscr I_\Gamma^{384},
 \qquad |\mathscr I_{\rm ext}|=855.
 \label{eq:frontal-inventario-masas-anchuras}
\end{equation}
CKM transporte les repères spectraux quark et PMNS les repères chargé et neutre
de rang trois ; leurs domaines, opérateurs et réfutateurs demeurent distincts.
Aucune valeur de l'inventaire ne sélectionne une voie, un coefficient, une valeur propre
ou un angle.
\end{theorem}

\begin{proof}
La forme normale de Smith produit
\(\operatorname{diag}(1,2,2,4)\) et un conoyau d'ordre \(16\) ; les bornes de
\(\Sigma_H\) fixent \(d_H=-34\), tandis que
\eqref{eq:frontal-mantisas-accion}--\eqref{eq:frontal-secciones-accion}
publient \(\hbar_a\) puis \(h_a\). La substitution de
\(\vartheta_{108}\) dans \(G_a(\vartheta)\) prouve
\eqref{eq:frontal-longitud-planck}. L'incidence des ensembles
\(H_{\rm act}\times\binom{H_{\rm act}}2\) et
\(\{1,\ldots,8\}\times\binom{H_{\rm act}}2\) produit \(90\) et \(120\).
Les douze secteurs dodécaphasiques et l'unique couture orientée de retour
produisent, au moyen du lecteur linéaire, le poids \(12{:}{-}1\) de la transition
hexade--octade. Le coefficient de \((1-q)^{-1}\), égal à \(1/24\), et la
minimalité diophantienne vérifient ensuite ce poids. Le recensement fini de l'atlas et de ses deux lecteurs
prouve les cardinaux de la branche de masse ; l'inversion cellulaire possède un unique
point fixe, \((5,5)\), qui réduit canoniquement la fibre électronique au rang
un. L'inventaire externe n'est composé qu'après avoir figé la fibre,
l'opérateur, le projecteur, la représentation, le schéma et l'échelle, de sorte qu'une
permutation de ses grandeurs laisse invariant tout le constructeur interne.
\end{proof}
% END HMT-MD-FRONTAL-MAXIMAL-SIGNATURES-20260904

\fi
